%% ============================================================
%%  Section: Empirical Strategy and Identification
%%  Depositor Discipline in Modern Banking (2026-04-03 draft)
%%
%%  NOTE: Written to cc_generated/v4/ due to session file limit.
%%  Move to latex/section_mds/identification/ before compiling.
%%
%%  Figures referenced:
%%    ../docs/figures/mkt_cap_large_banks_20260403.png
%%    ../docs/figures/cds_spread_large_banks_20260403.png
%% ============================================================

\section{Empirical Strategy and Identification}
\label{sec:identification}

\subsection{Difference-in-Differences Design}
\label{sec:identification:did}

Our identification strategy exploits cross-sectional variation in the size of
the CECL Day-One retained-earnings adjustment, scaled by pre-adoption equity, as
a bank-level treatment intensity. We estimate difference-in-differences
specifications of the form
\begin{equation}
  y_{it} = \alpha_i + \gamma_t + \beta_1 \, \text{Post}_{it} \times \text{CECL}_i
  + \mathbf{X}_{it-1}'\boldsymbol{\delta} + \varepsilon_{it},
  \label{eq:did}
\end{equation}
where $y_{it}$ is the outcome for bank $i$ in quarter $t$, $\alpha_i$ and
$\gamma_t$ are bank and calendar-quarter fixed effects, $\text{Post}_{it}$
equals one for all quarters at or after bank $i$'s CECL adoption date, and
$\text{CECL}_i$ is either the continuous CECL-to-equity adjustment measured at
adoption or a binary indicator for the top quartile of that distribution
($\text{High CECL}_i$). The coefficient $\beta_1$ captures the differential
change in outcomes for higher-exposure banks after adoption, relative to
lower-exposure banks, net of bank-level trends absorbed by the fixed effects.
$\mathbf{X}_{it-1}$ includes lagged log assets, lagged equity-to-assets, and
lagged interest expense-to-assets, all measured one quarter prior to $t$.

The estimating sample consists of community banks with total assets below
\$10~billion as of 2022Q4, observed from 2016Q1 through 2025Q4. We restrict
attention to banks whose CECL adoption quarter falls within a six-quarter window
centered on 2023Q1 (from 2022Q3 through 2023Q4), the standard rollout window for
non-SEC-filing institutions. Community banks outside this window---those with no
detectable adoption date---are excluded, as are the roughly 200 large SEC-filing
banks that adopted in 2020Q1, for reasons explained in
Section~\ref{sec:identification:signal} below.

\subsection{Identification: CECL as a Predetermined Information Shock}
\label{sec:identification:predetermination}

The core identifying assumption is that the CECL adjustment is predetermined
relative to the post-adoption deposit and financial outcomes we study. This
assumption has three distinct components.

\paragraph{Content predetermination.}
The CECL adjustment is computed as the difference between the new expected credit
loss allowance under ASC~326 and the legacy incurred-loss reserve, applied to
the loan portfolio as it stood on the adoption date. It reflects the accumulated
stock of credit risk embedded in existing loans at a fixed point in time, not
management's forward-looking response to post-adoption conditions. Crucially, it
is disclosed simultaneously for all adopting banks on their first quarterly
filing under CECL, meaning depositors receive the information in a single
cross-sectional release rather than through gradual revelation. The size of the
adjustment is therefore determined by loan portfolio characteristics frozen before
the disclosure, not by the deposit outcomes that follow.

\paragraph{Timing predetermination.}
Community banks chose their fiscal-year CECL adoption dates years before the
disclosure, subject to regulatory deadlines set by the FASB and federal banking
agencies. A bank that adopted in 2023Q1 had committed to that timeline no later
than 2019; a bank that adopted in 2023Q4 had similarly pre-committed. Neither
bank could have adjusted its adoption timing in response to what the adjustment
would reveal about its credit risk, as the adjustment itself was not calculated
until the adoption quarter. This eliminates the possibility that banks with
larger expected adjustments delayed adoption to avoid adverse depositor responses.

\paragraph{Capital phase-in predetermination.}
Under the 2020 interim final rule, banks electing the three-year regulatory
capital transition were permitted to phase in the retained-earnings charge to
their regulatory capital ratios over 2022--2024. The election was required to be
made in advance of adoption. Banks that elected the phase-in
($\text{Phase\_In}_i = 1$) faced the same informational disclosure to depositors
as non-electing banks---the full CECL adjustment appears in the financial
statements regardless of the phase-in election---but experienced a more gradual
mechanical capital impact. This separation between information disclosure and
capital impact provides an additional check on the mechanism: if depositor
responses are driven by information rather than capital, they should be present
regardless of phase-in status.

\subsection{Signal Validation: Market and CDS Responses in the SEC-Filer Cohort}
\label{sec:identification:signal}

A necessary condition for our interpretation is that the CECL adjustment
contains genuine new information about bank credit quality---information that
sophisticated outside observers could not have extracted from prior
disclosures. If CECL were merely a relabeling of already-known reserve
positions, no rational depositor or investor should respond to it. We test this
condition using the cohort of large, publicly traded banks that adopted CECL in
2020Q1 under SEC-filer requirements.

\paragraph{Why the 2020 cohort, and why only for validation.}
SEC-filing banks (primarily banks with assets above \$10~billion) were required
to adopt CECL in their first fiscal year beginning after December~15, 2019, which
for calendar-year filers means 2020Q1. This cohort adopted on a known, uniform
date---unlike community banks, whose adoption was staggered across 2022Q4 through
2023Q4---so the event time is sharply defined. However, 2020Q1 coincides
precisely with the onset of the COVID-19 pandemic. The aggregate shock to bank
credit risk in that quarter is so large, and so correlated with CECL adoption
timing, that any attempt to use the 2020 cohort to estimate depositor responses
would be confounded by COVID-driven deposit flows. We therefore restrict our
main analysis to community banks adopting in 2023, where no comparable aggregate
shock coincides with the disclosure. We use the 2020 SEC-filer cohort only to
establish that the CECL adjustment carries genuine information content that moves
sophisticated financial market participants---a validation exercise for the
instrument, not an attempt to identify depositor discipline.

\paragraph{Equity market response.}
Figure~\ref{fig:es_large_banks_market_cds} Panel A plots event-study estimates of the
differential change in normalized equity market capitalization for SEC-filer banks
in the nine months before and after their CECL adoption date, using the
continuous CECL-to-equity ratio as the treatment intensity. All specifications
include bank and calendar-month fixed effects and are estimated on the subsample
of adopters through 2021 to maintain distance from the COVID period.\footnote{The
pre-adoption window begins at $m = -9$ and the reference period is $m = -1$.
Banks with earlier CECL adoption dates (2019) are included; the COVID-era 2020
cohort is partially included with the caveat that month fixed effects absorb
common pandemic variation.}

The pre-adoption coefficients are flat and statistically indistinguishable from
zero across all nine lead months, confirming that equity markets did not
differentially price high-CECL banks before the disclosure. At adoption
(month~0) the coefficient drops sharply to $-0.947$ ($p = 0.006$), indicating
that a one-percentage-point higher CECL-to-equity ratio is associated with a
$0.95$-unit decline in normalized market capitalization in the adoption month
itself. The effect deepens through months~1 and~2, reaching a trough of
$-2.771$ ($p < 0.001$) at month~2---implying that a ten-percentage-point higher
CECL adjustment is associated with roughly a $28$~percent lower market
capitalization relative to adoption-month levels, after removing common monthly
shocks and bank fixed effects. The effect then partially reverts over months~3
through~9 but remains negative and economically substantial through month~6. This
pattern is inconsistent with the CECL adjustment being merely a mechanical
accounting restatement: a purely cosmetic change to reserve classifications would
carry no information about underlying credit quality and should produce no
market cap response.

\paragraph{Credit default swap response.}
Figure~\ref{fig:es_large_banks_market_cds} Panel B repeats the event-study analysis using
monthly CDS spreads for the subset of SEC-filing banks with active CDS contracts
and a complete nine-month symmetric window around adoption (103 banks). The CDS
market provides a complementary signal: while equity values reflect the residual
claim on the bank's asset portfolio, CDS spreads directly price the probability
of default and are traded by sophisticated institutional counterparties
specifically because they serve as credit risk monitors.

The pre-adoption CDS coefficients are flat and cluster tightly around zero
through all nine lead months---even more precisely than the equity
estimates---reflecting the concentrated sophistication of CDS market
participants. At adoption, CDS spreads jump immediately and rise persistently
through month~3. The most precisely estimated post-adoption coefficient occurs at
month~4: $+190.8$~basis points per unit of CECL-to-equity ($p = 0.047$),
implying that a ten-percentage-point higher CECL adjustment is associated with
approximately 19~basis points wider CDS spreads four months after adoption.
Months~2 ($+164$~bps, $p = 0.090$) and~5 ($+152$~bps, $p = 0.056$) are
marginally significant. The response is noisier than the equity response, as
expected given the smaller CDS sample, but the directional pattern is
unambiguous: the CDS market persistently prices higher default risk at
high-CECL banks following the disclosure.

\paragraph{Joint interpretation.}
The equity and CDS results together establish two things. First, the CECL
adjustment is informative: sophisticated financial market participants with every
incentive to have already priced bank credit risk revised their valuations
sharply and persistently at the moment of CECL adoption. The flat pre-trends in
both markets rule out the possibility that the adjustment was anticipated. Second,
the information content is specifically about credit risk, not about other
mechanical consequences of CECL adoption such as regulatory capital changes: CDS
spreads price default probability directly, and their response cannot be
attributed to capital-ratio mechanics. These two properties---genuine information
content, unanticipated timing---are exactly what is required for the CECL
disclosure to serve as a valid shock in our community-bank depositor discipline
design.